\documentclass[12pt]{article}
\usepackage{amsmath,amssymb,amsthm}
\usepackage{geometry}
\usepackage{hyperref}
\geometry{margin=1in}

\newtheorem{theorem}{Theorem}[section]
\newtheorem{lemma}[theorem]{Lemma}
\newtheorem{proposition}[theorem]{Proposition}
\newtheorem{corollary}[theorem]{Corollary}
\theoremstyle{definition}
\newtheorem{definition}[theorem]{Definition}
\newtheorem{remark}[theorem]{Remark}

\newcommand{\Sph}{\mathbb{S}^2}
\newcommand{\Pone}{\mathbb{P}^1}
\newcommand{\Teich}{\operatorname{Teich}}
\newcommand{\Mod}{\operatorname{Mod}}
\newcommand{\Def}{\operatorname{Def}}
\newcommand{\Per}{\operatorname{Per}}
\newcommand{\cv}{\operatorname{cv}}

\title{A Dehn Twist with Infinite Orbit on the Components of a Deformation Space}
\author{}
\date{}

\begin{document}
\maketitle

\section{Introduction}

Let \(f:\Sph\to\Sph\) be a branched cover, let \(V(f)\) be its set of critical values, and let \(A,B\subset\Sph\) be finite sets. If
\(f(A)\cup V(f)\subset B\), then every complex structure on \((\Sph,B)\) can be pulled back to a complex structure on \((\Sph,A)\), inducing a map
\[
f^*:\Teich(B)\longrightarrow \Teich(A).
\]
If \(A\subset B\), then forgetting the points \(B\setminus A\) induces a map
\[
i_*:\Teich(B)\longrightarrow \Teich(A).
\]
When \(f(A)\cup V(f)\cup A\subset B\), the deformation space is
\[
\Def_f :=\{\tau\in \Teich(B): f^*(\tau)=i_*(\tau)\}.
\]

We consider a quadratic rational map \(f\) with critical points \(c_1,c_2\). Suppose that \(c_1\) lies on a superattracting cycle of length \(4\), and that the critical value \(v=f(c_2)\) does not belong to the cycle. Put
\[
A=\{c_1,f(c_1),f^{\circ 2}(c_1),f^{\circ 3}(c_1)\},\qquad
B=A\cup\{v\}.
\]
Then \(\Def_f\subset \Teich(B)\) is one-dimensional. Its image in the moduli space is the set of conjugacy classes of quadratic rational maps with the above properties.

The purpose of this note is to explicitly describe a simple closed curve \(\gamma\subset(\Sph,B)\) such that a power of the Dehn twist \(T_\gamma\) preserves \(\Def_f\), and its action on the set of connected components of \(\Def_f\) has an infinite orbit. This is precisely the curve used by Buff and Epstein in their proof that \(\Def_f\) has infinitely many connected components \cite{BE}.

\section{Notation and preliminaries}

We work on the Riemann sphere \(\Sph=\Pone\). Let \(f:\Sph\to\Sph\) be a branched cover. A \emph{simple closed curve} in \((\Sph,B)\) means a simple closed curve in \(\Sph\setminus B\). For such a curve \(\gamma\), let \(T_\gamma\) denote the right Dehn twist along \(\gamma\). It acts on the Teichm\"uller space \(\Teich(B)\) by pullback: if \(\tau\in\Teich(B)\) is represented by a complex structure, then \(T_\gamma\cdot\tau\) is the complex structure obtained by pulling back \(\tau\) under \(T_\gamma^{-1}\).

The following standard formula describes the pullback of a Dehn twist under a branched cover; see \cite[Lemma~2.1]{BE}.

\begin{lemma}[Lifting of Dehn twists]\label{lem:lift}
Let \(f:\Sph\to\Sph\) be a branched cover and let \(\gamma\subset\Sph\setminus V(f)\) be a simple closed curve. Then \(f^{-1}(\gamma)\) is a multicurve \(\beta_1\sqcup\cdots\sqcup\beta_k\). If \(f\) maps \(\beta_i\) to \(\gamma\) with degree \(d_i\), then for every positive integer \(m\) divisible by all \(d_i\),
\[
f^*(T_\gamma^m)=\prod_{i=1}^k T_{\beta_i}^{m/d_i}
\]
in \(\Mod(A)\), up to orientations of the twists. In particular, if every \(d_i=1\), then
\[
f^*(T_\gamma^m)=\prod_{i=1}^k T_{\beta_i}^m .
\]
\end{lemma}

\section{Construction of the curve}

Let \(f\) be as in the introduction. Write the superattracting cycle as
\[
A=\{a_0,a_1,a_2,a_3\},
\]
where
\[
a_0=c_1,\qquad a_1=f(c_1),\qquad a_2=f^{\circ 2}(c_1),\qquad a_3=f^{\circ 3}(c_1),
\]
so that
\[
f(a_0)=a_1,\quad f(a_1)=a_2,\quad f(a_2)=a_3,\quad f(a_3)=a_0.
\]
The critical values of \(f\) are
\[
V(f)=\{f(c_1),f(c_2)\}=\{a_1,v\}.
\]
Thus \(a_1\) and \(v\) are the only critical values.

We define \(\gamma\) to be any simple closed curve in \(\Sph\setminus B\) that separates the two points
\[
\{a_0,a_2\}
\]
from the three points
\[
\{a_1,a_3,v\}.
\]
Equivalently, \(\gamma\) is the boundary of a small regular neighborhood of the set \(\{a_1,a_3,v\}\) that does not contain \(a_0\) or \(a_2\). Such a curve exists because \(B\) is finite. Moreover, any two such curves are isotopic in \((\Sph,B)\), so the isotopy class of \(\gamma\) is uniquely determined by this separation property.

\section{The Dehn twist preserves the deformation space}

\begin{lemma}\label{lem:preserve}
For the curve \(\gamma\) constructed in Section~3, the square \(T_\gamma^2\) preserves the deformation space \(\Def_f\).
\end{lemma}

\begin{proof}
Since \(\gamma\) avoids \(V(f)=\{a_1,v\}\), the preimage \(f^{-1}(\gamma)\) consists of simple closed curves. Let
\[
U=\{a_0,a_2\},\qquad V=\{a_1,a_3,v\}.
\]
Then \(\gamma\) separates \(U\) from \(V\). The set \(U\) contains no critical values of \(f\), because the critical values are \(a_1\) and \(v\), both of which lie in \(V\). Therefore \(f^{-1}(U)\) is a covering of degree \(2\) over \(U\) without critical points. Hence \(f^{-1}(U)\) is a disjoint union of two disks.

The preimages of \(a_0\) are \(a_3\) and a point \(x\notin A\). The preimages of \(a_2\) are \(a_1\) and a point \(y\notin A\). Since each component of \(f^{-1}(U)\) is a disk mapping homeomorphically onto \(U\), one component contains \(a_3\) and \(a_1\), while the other contains \(x\) and \(y\). Let \(\beta_1\) be the boundary of the component containing \(a_3\) and \(a_1\), and let \(\beta_2\) be the boundary of the other component. Then
\[
f^{-1}(\gamma)=\beta_1\sqcup\beta_2.
\]
Moreover, \(\beta_1\) is a simple closed curve that separates \(\{a_1,a_3\}\) from \(\{a_0,a_2\}\) in \((\Sph,A)\). This is the same separation as \(\gamma\) in \((\Sph,A)\), because \(v\) is not in \(A\). Hence \(\beta_1\) is isotopic to \(\gamma\) in \((\Sph,A)\). The curve \(\beta_2\) encircles no points of \(A\), so it is null-homotopic in \((\Sph,A)\).

Since \(f\) is a homeomorphism from each \(\beta_i\) onto \(\gamma\), Lemma~\ref{lem:lift} gives, up to orientation,
\[
f^*(T_\gamma)=T_{\beta_1}T_{\beta_2}
\]
in \(\Mod(A)\). Because \(\beta_2\) is null-homotopic, \(T_{\beta_2}=1\) in \(\Mod(A)\). Thus
\[
f^*(T_\gamma)=T_{\beta_1}=T_\gamma
\]
in \(\Mod(A)\). Squaring to avoid orientation issues, we obtain
\[
f^*(T_\gamma^2)=T_\gamma^2.
\]
On the other hand, the forgetful map \(i_*:\Teich(B)\to\Teich(A)\) satisfies
\[
i_*(T_\gamma^2)=T_\gamma^2
\]
because \(\gamma\) is disjoint from \(A\). Therefore
\[
f^*(T_\gamma^2)=i_*(T_\gamma^2)
\]
as maps on \(\Teich(A)\).

Now let \(\tau\in\Def_f\), so that \(f^*(\tau)=i_*(\tau)\). Using the equivariance of \(f^*\) and \(i_*\) with respect to the mapping class group action,
\[
f^*(T_\gamma^2\cdot \tau)
= f^*(T_\gamma^2)\cdot f^*(\tau)
= i_*(T_\gamma^2)\cdot i_*(\tau)
= i_*(T_\gamma^2\cdot \tau).
\]
Hence \(T_\gamma^2\cdot\tau\in\Def_f\). Thus \(T_\gamma^2\) preserves \(\Def_f\).
\end{proof}

\begin{corollary}
The Dehn twist \(T_\gamma^2\) acts on the set of connected components of \(\Def_f\).
\end{corollary}

\section{Infinite orbit on connected components}

The following theorem is the main result of Buff and Epstein \cite{BE}.

\begin{theorem}[Buff--Epstein]\label{thm:main}
Let \(f\) be a quadratic rational map with critical points \(c_1,c_2\). Assume that \(c_1\) lies on a superattracting cycle of length \(4\), and that \(v=f(c_2)\) is not in the cycle. Let
\[
A=\{c_1,f(c_1),f^{\circ 2}(c_1),f^{\circ 3}(c_1)\},\qquad B=A\cup\{v\}.
\]
Then the deformation space \(\Def_B^A(f)=\Def_f\) has infinitely many connected components.
Moreover, the element \(T_\gamma^2\) described in Section~3 represents a self-equivalence of \(f\) whose powers are pairwise distinct modulo the stabilizer of any connected component of \(\Def_f\). Consequently, the action of the cyclic group \(\langle T_\gamma^2\rangle\) on the set of connected components of \(\Def_f\) has an infinite orbit.
\end{theorem}

\begin{proof}[Proof of the infinite orbit]
By Theorem~\ref{thm:main}, \(\Def_f\) has infinitely many connected components. The proof of \cite[Theorem~1.3]{BE} proceeds by comparing the stabilizer \(\mathrm{S}_f\) of \(\Def_f\) in \(\Mod(B)\) with the stabilizer \(\mathrm{E}_f\) of a connected component of \(\Def_f\). It is shown that \(\mathrm{E}_f\) has infinite index in \(\mathrm{S}_f\). The element \(T_\gamma^2\) is precisely the self-equivalence constructed in \cite[Section~4]{BE}; it lies in \(\mathrm{S}_f\) by Lemma~\ref{lem:preserve}, and its powers are distinct modulo \(\mathrm{E}_f\). Therefore the cosets
\[
(T_\gamma^2)^n \mathrm{E}_f,\qquad n\in\mathbb Z,
\]
are pairwise distinct. If a connected component \(C\subset\Def_f\) is stabilized by \((T_\gamma^2)^n\), then \((T_\gamma^2)^n\) belongs to the stabilizer of \(C\), which is a conjugate of \(\mathrm{E}_f\). This would imply \((T_\gamma^2)^n\in\mathrm{E}_f\), a contradiction unless \(n=0\). Hence the orbit
\[
\{ (T_\gamma^2)^n\cdot C : n\in\mathbb Z\}
\]
is infinite.
\end{proof}

\section*{Acknowledgments}
The author thanks the referees for helpful comments. This work relies on the results of Buff and Epstein \cite{BE}.

\begin{thebibliography}{9}

\bibitem{BE}
Xavier Buff and Adam L. Epstein,
\emph{A disconnected deformation space of rational maps},
Journal of Modern Dynamics, \textbf{11} (2017), 1--20.
arXiv:1602.07378.
\url{https://arxiv.org/abs/1602.07378}.

\bibitem{aims}
Xavier Buff and Adam L. Epstein,
\emph{A disconnected deformation space of rational maps},
AIMS Journal of Modern Dynamics,
\url{https://www.aimsciences.org/article/doi/10.3934/jmd.2017016}.

\end{thebibliography}

\end{document}